% Options for packages loaded elsewhere
\PassOptionsToPackage{unicode}{hyperref}
\PassOptionsToPackage{hyphens}{url}
%
\documentclass[
  11pt,
]{article}
\usepackage{amsmath,amssymb}
\usepackage{iftex}
\ifPDFTeX
  \usepackage[T1]{fontenc}
  \usepackage[utf8]{inputenc}
  \usepackage{textcomp} % provide euro and other symbols
\else % if luatex or xetex
  \usepackage{unicode-math} % this also loads fontspec
  \defaultfontfeatures{Scale=MatchLowercase}
  \defaultfontfeatures[\rmfamily]{Ligatures=TeX,Scale=1}
\fi
\usepackage{lmodern}
\ifPDFTeX\else
  % xetex/luatex font selection
\fi
% Use upquote if available, for straight quotes in verbatim environments
\IfFileExists{upquote.sty}{\usepackage{upquote}}{}
\IfFileExists{microtype.sty}{% use microtype if available
  \usepackage[]{microtype}
  \UseMicrotypeSet[protrusion]{basicmath} % disable protrusion for tt fonts
}{}
\makeatletter
\@ifundefined{KOMAClassName}{% if non-KOMA class
  \IfFileExists{parskip.sty}{%
    \usepackage{parskip}
  }{% else
    \setlength{\parindent}{0pt}
    \setlength{\parskip}{6pt plus 2pt minus 1pt}}
}{% if KOMA class
  \KOMAoptions{parskip=half}}
\makeatother
\usepackage{xcolor}
\usepackage[margin=1in]{geometry}
\usepackage{longtable,booktabs,array}
\usepackage{calc} % for calculating minipage widths
% Correct order of tables after \paragraph or \subparagraph
\usepackage{etoolbox}
\makeatletter
\patchcmd\longtable{\par}{\if@noskipsec\mbox{}\fi\par}{}{}
\makeatother
% Allow footnotes in longtable head/foot
\IfFileExists{footnotehyper.sty}{\usepackage{footnotehyper}}{\usepackage{footnote}}
\makesavenoteenv{longtable}
\setlength{\emergencystretch}{3em} % prevent overfull lines
\providecommand{\tightlist}{%
  \setlength{\itemsep}{0pt}\setlength{\parskip}{0pt}}
\setcounter{secnumdepth}{-\maxdimen} % remove section numbering
\ifLuaTeX
  \usepackage{selnolig}  % disable illegal ligatures
\fi
\usepackage{bookmark}
\IfFileExists{xurl.sty}{\usepackage{xurl}}{} % add URL line breaks if available
\urlstyle{same}
\hypersetup{
  pdftitle={Incompleteness-Secured Commitments: a Dual-Layer Post-Quantum Framework with Reproducibility Guards, Wrapper Portability Probes, and an ML-KEM-768 Reference Prototype},
  pdfauthor={Kevin Henry Miller; Q-Bond Network DeSCI DAO, LLC; kevin@qbondnetwork.com},
  hidelinks,
  pdfcreator={LaTeX via pandoc}}

\title{Incompleteness-Secured Commitments: a Dual-Layer Post-Quantum
Framework with Reproducibility Guards, Wrapper Portability Probes, and
an ML-KEM-768 Reference Prototype}
\author{Kevin Henry Miller \and Q-Bond Network DeSCI DAO,
LLC \and kevin@qbondnetwork.com}
\date{March 2026}

\begin{document}
\maketitle

\section{Abstract}\label{abstract}

We describe a commitment framework that combines an ordinary
post-quantum cryptographic body with a proof-carrying logical annotation
layer. The cryptographic body is built around ML-KEM-768 as a
key-encapsulation mechanism plus a one-bit masking layer and a
transcript hash. The logical layer attaches a canonical arithmetic
formula encoding the opening relation and uses Gödel numbering only to
define a proof-carrying meta-security game; it does not replace ordinary
computational security reductions. In this formulation, hiding and
binding remain conventional and are argued at the concrete cryptographic
layer, while the additional claim is narrower: a successful
\emph{certified} double opening must also be accompanied by a formal
proof inside a fixed theory that both openings verify. Under standard
representability assumptions, such a certified attack would amount to a
proof of a false arithmetic statement whenever the concrete
cryptographic layer has not already failed.

The paper is deliberately conservative. It does not claim
information-theoretic security, production readiness, constant-time
behavior for the Python orchestration layer, or a complete IND-CPA,
IND-CCA, or UC proof of the full construction. Instead, it focuses on
three practical deliverables: an arXiv-friendly manuscript, a public
reference repository, and measured reproducibility artifacts. We report
an archived last-good baseline line, a second archived continuity line,
a current environment rerun on the legacy ML-KEM alpha route, and a
wrapper-portability audit for keeper-ML-KEM and libfips203 skeleton
adapters. The result is a clearer separation between what has been
formalized, what has been implemented, and what remains future work.

\section{1. Introduction}\label{introduction}

Modern post-quantum cryptography is built on computational assumptions,
and that is the correct baseline. The natural question explored here is
not whether Gödel's incompleteness theorem can magically strengthen
ciphertexts by itself, but whether a \emph{proof-carrying} variant of a
standard commitment system can cleanly separate two layers of failure.
The first layer is conventional cryptographic failure: for example, a
hiding or binding break under an implementation or assumption-level
attack. The second layer is logical certification of that failure: an
adversary not only outputs two openings to the same commitment, but also
outputs a formal derivation inside a fixed theory certifying that both
openings verify.

That second requirement is where incompleteness enters. The verifier for
the commitment is decidable. Therefore the statement ``this commitment
admits two distinct valid openings'' can be represented in arithmetic.
If the concrete cryptographic layer is still behaving as intended, then
the represented sentence is false. A proof-carrying adversary is
therefore trying to package a cryptographic event together with a formal
proof of a false statement. The point of the framework is not to replace
standard cryptographic analysis. The point is to isolate a narrow and
explicit interface between conventional cryptography and proof theory.

The manuscript is written against a concrete practical problem: earlier
drafts mixed cryptographic claims, package claims, and broader
speculative material too aggressively. This version deliberately
tightens scope. Older package-specific mentions have been removed. The
software is described as an experimental reference prototype.
Secret-dependent work is described as delegated to the KEM backend
rather than to Python glue. A manuscript claim-guard script has been
added so that high-risk overclaim language, unsupported deployment
language, or category-driven positioning language cannot silently creep
back into the paper.

\section{2. Contributions and scope
discipline}\label{contributions-and-scope-discipline}

The paper makes six bounded contributions.

\begin{enumerate}
\def\labelenumi{\arabic{enumi}.}
\tightlist
\item
  It defines a dual-layer commitment framework in which the concrete
  commitment body is standard and the logical layer is explicitly
  meta-security only.
\item
  It states a proof-carrying adversary model and a corresponding notion
  of certified double-opening failure.
\item
  It implements a reference commitment prototype using ML-KEM-768 as a
  \emph{key-encapsulation mechanism}, not as an improvised public-key
  encryption primitive.
\item
  It publishes a public repository with scripts for backend probing,
  smoke tests, proof-bundle checks, benchmark runs, and a manuscript
  claim guard.
\item
  It records archived benchmark lines for continuity and adds a current
  rerun line from the present workspace.
\item
  It performs a wrapper portability audit for keeper-ML-KEM and
  libfips203 skeleton adapters and records the exact observed failure
  modes.
\end{enumerate}

The scope discipline matters. The wider Q-Bond / TICE / Floquet /
κ-meter corpus includes many ambitious materials spanning quantum
control, metrology, and information-curvature ideas, including broad
technical notebooks and patent-style drafts. That material is real
context for the project as a whole, but importing those claims directly
into a cryptography paper would weaken rather than strengthen the
present manuscript. This draft therefore borrows only one lesson from
that larger corpus: the emphasis on packaging code, measurements, and
reproducibility artifacts together with the text. The concrete
cryptography claims stand on their own and are intentionally narrower
than many of the adjacent research narratives.

\section{3. Model and construction}\label{model-and-construction}

\subsection{3.1 Concrete commitment
body}\label{concrete-commitment-body}

The reference commitment body is

\[
B=(ct_{\mathsf{kem}}, d, 
ho),
\]

where \(ct_{\mathsf{kem}}\) is an ML-KEM-768 ciphertext, \$ ho\$ is a
32-byte public nonce, and \(d\in\{0,1\}\) is a masked bit carried on the
wire as one byte. The full commitment is then

\[
C=(B,c_G),
\]

where \(c_G\) is a SHA3-256 transcript hash over a canonical
Gödel-encoded arithmetic statement and the opening randomness.

Operationally, the prototype proceeds as follows.

\begin{enumerate}
\def\labelenumi{\arabic{enumi}.}
\tightlist
\item
  Generate an ML-KEM-768 keypair \((pk,sk)\).
\item
  Run \(\mathsf{Encaps}(pk)\) to obtain a shared secret \(ss\) and a KEM
  ciphertext \(ct_{\mathsf{kem}}\).
\item
  Derive a public nonce \$ ho\$ from explicit randomness and the
  ciphertext.
\item
  Derive a one-bit mask from \((ss,
  ho)\) and publish the masked bit \(d=m\oplus b\).
\item
  Build a canonical arithmetic formula \(\phi\) describing the opening
  relation for the current transcript.
\item
  Gödel-encode \(\phi\) and compute \[
  c_G = \mathrm{SHA3   ext{-}256}(\mathrm{enc}(\phi)\parallel B\parallel r).
  \]
\item
  Publish \(C=(B,c_G)\) and retain the opening witness \(w=(m,r,\phi)\).
\end{enumerate}

Opening uses decapsulation to recompute the mask and recover the bit.
Verification is public and checks that the canonical formula recomputed
from \((B,m,r)\) matches the supplied formula and that the transcript
digest recomputes exactly.

\subsection{3.2 What remains standard}\label{what-remains-standard}

The cryptographic layer remains conventional. Hiding depends on the
encapsulation-and-mask path. Binding depends on the transcript hash, the
canonicalization routine, and the underlying backend behaving
consistently. Nothing in the paper claims that incompleteness upgrades
these ordinary guarantees into unconditional security. The construction
remains a normal cryptographic object with a secondary logical game
attached to it.

\subsection{3.3 Proof-carrying adversary
model}\label{proof-carrying-adversary-model}

Let \(\mathsf{Bad}(C,w_0,w_1)\) denote the arithmetic sentence stating
that commitment \(C\) admits two distinct valid openings \(w_0\) and
\(w_1\). A proof-carrying adversary succeeds only if it outputs not
merely the commitment and two openings, but also a derivation \(\pi\) in
a fixed theory \(T\) that certifies \(\mathsf{Bad}(C,w_0,w_1)\).

The role of Gödel numbering here is narrow. It is a bridge from the
verification relation into arithmetic syntax so that the proof-carrying
game is well-defined. It is not a claim that every practical attack on
the implementation must route through theorem proving. The only added
security statement is about a stronger attacker model that wants a
certified break.

\section{4. Informal security
discussion}\label{informal-security-discussion}

The concrete hiding story is the familiar one: before opening, an
observer sees a KEM ciphertext, a nonce, a one-byte masked bit, and a
transcript digest. If the shared secret derived by encapsulation remains
hidden and the one-bit mask is derived from that secret and the public
nonce, then the published masked bit does not reveal the committed bit.
The concrete binding story is equally familiar: two successful openings
for the same commitment should either indicate a collision in the
transcript hash or a defect in the canonicalization routine or backend
behavior.

The meta-security statement is weaker and more specialized. If the
concrete layer has not already failed, then a sentence asserting the
existence of a valid double opening is false. A consistent and
sufficiently expressive theory should therefore not prove it. In that
sense, incompleteness does not make the commitment body
cryptographically stronger by itself; it makes \emph{certified
falsehood} harder in the proof-carrying game.

This distinction is worth repeating because it prevents the common
failure mode of overstatement. The scheme is not above all
computation-based systems in the ordinary cryptographic sense. It is a
standard post-quantum commitment architecture with an additional,
logically flavored certificate layer.

\section{5. Reference implementation}\label{reference-implementation}

\subsection{5.1 Why the implementation path
matters}\label{why-the-implementation-path-matters}

The current reference implementation uses ML-KEM-768 in the way the
standard intends: as a \emph{key-encapsulation mechanism}. The paper
therefore avoids descriptions such as ``encrypt a bit with Kyber.''
Instead, the bit is masked using material derived from the shared
secret, and the commitment transcript is bound by a SHA3-256 digest.

\subsection{5.2 Wire format}\label{wire-format}

The current wire format is shown below.

\begin{longtable}[]{@{}
  >{\raggedright\arraybackslash}p{(\columnwidth - 4\tabcolsep) * \real{0.3000}}
  >{\raggedleft\arraybackslash}p{(\columnwidth - 4\tabcolsep) * \real{0.4000}}
  >{\raggedright\arraybackslash}p{(\columnwidth - 4\tabcolsep) * \real{0.3000}}@{}}
\toprule\noalign{}
\begin{minipage}[b]{\linewidth}\raggedright
Component
\end{minipage} & \begin{minipage}[b]{\linewidth}\raggedleft
Size
\end{minipage} & \begin{minipage}[b]{\linewidth}\raggedright
Comment
\end{minipage} \\
\midrule\noalign{}
\endhead
\bottomrule\noalign{}
\endlastfoot
ML-KEM-768 ciphertext & 1088 B & KEM ciphertext \\
Public nonce \$ & & \\
ho\$ & 32 B & nonce derived from randomness and ciphertext \\
Masked bit \(d\) & 1 B & one-byte wire encoding \\
Transcript digest \(c_G\) & 32 B & SHA3-256 over canonical statement and
opening data \\
\textbf{Total} & \textbf{1153 B} & current commitment size \\
\end{longtable}

That 1153-byte size remains stable across the archived benchmark lines
and the current workspace rerun line.

\subsection{5.3 Public repository
contents}\label{public-repository-contents}

The public repository accompanying this paper includes the following
items.

\begin{itemize}
\tightlist
\item
  \texttt{src/gce/engine.py}: commit, verify, and open logic.
\item
  \texttt{src/gce/kem\_backends.py}: backend selectors for keeper-style,
  legacy-alpha, libfips203, and simulated lanes.
\item
  \texttt{src/gce/kem\_wrappers/keeper\_mlkem\_wrapper.py}: skeleton
  adapter for a Keeper-style ML-KEM package.
\item
  \texttt{src/gce/kem\_wrappers/libfips203\_wrapper.py}: skeleton
  adapter for Python \texttt{fips203} / \texttt{libfips203} style
  bindings.
\item
  \texttt{scripts/00\_backend\_probe.py}: report backend metadata and
  wire-format sizes.
\item
  \texttt{scripts/01\_smoke\_test.py}: minimal round-trip check.
\item
  \texttt{scripts/02\_proof\_bundle.py}: round-trip, wrong-opening
  rejection, tamper-detection, and Gödel-number uniqueness checks.
\item
  \texttt{scripts/03\_benchmark.py}: timing benchmark.
\item
  \texttt{scripts/04\_cross\_backend\_validate.py}: compare multiple
  backend lanes.
\item
  \texttt{scripts/05\_probe\_wrapper\_status.py}: probe the skeleton
  wrappers and record exact availability or failure modes.
\item
  \texttt{scripts/check\_claims.py}: fail if banned phrases or missing
  disclaimer language appear in the manuscript.
\end{itemize}

The claim guard is not cosmetic. It exists because earlier drafts
drifted too easily into language that was not supported by the
implementation or by the threat model.

\section{6. Measured results}\label{measured-results}

\subsection{6.1 Archived continuity
lines}\label{archived-continuity-lines}

The paper preserves two archived measurement lines because the user
requested that continuity record explicitly. Those lines describe prior
local validated runs at the same 1153-byte commitment size.

\begin{longtable}[]{@{}
  >{\raggedright\arraybackslash}p{(\columnwidth - 10\tabcolsep) * \real{0.1304}}
  >{\raggedleft\arraybackslash}p{(\columnwidth - 10\tabcolsep) * \real{0.1739}}
  >{\raggedleft\arraybackslash}p{(\columnwidth - 10\tabcolsep) * \real{0.1739}}
  >{\raggedleft\arraybackslash}p{(\columnwidth - 10\tabcolsep) * \real{0.1739}}
  >{\raggedleft\arraybackslash}p{(\columnwidth - 10\tabcolsep) * \real{0.1739}}
  >{\raggedleft\arraybackslash}p{(\columnwidth - 10\tabcolsep) * \real{0.1739}}@{}}
\toprule\noalign{}
\begin{minipage}[b]{\linewidth}\raggedright
Archived line
\end{minipage} & \begin{minipage}[b]{\linewidth}\raggedleft
Keygen ms
\end{minipage} & \begin{minipage}[b]{\linewidth}\raggedleft
Commit ms
\end{minipage} & \begin{minipage}[b]{\linewidth}\raggedleft
Verify ms
\end{minipage} & \begin{minipage}[b]{\linewidth}\raggedleft
Open ms
\end{minipage} & \begin{minipage}[b]{\linewidth}\raggedleft
Size B
\end{minipage} \\
\midrule\noalign{}
\endhead
\bottomrule\noalign{}
\endlastfoot
Archived last-good baseline & 0.566 & 0.861 & 0.023 & 0.902 & 1153 \\
Archived second continuity line & 0.580 & 0.870 & 0.025 & 0.910 &
1153 \\
\end{longtable}

These are retained as provenance and continuity data, not as newly rerun
measurements from the present workspace.

\subsection{6.2 Current workspace rerun}\label{current-workspace-rerun}

The current workspace rerun is different. In this environment, the
legacy ML-KEM alpha path remains runnable and produced the following
benchmark line over 20 trials.

\begin{longtable}[]{@{}
  >{\raggedright\arraybackslash}p{(\columnwidth - 12\tabcolsep) * \real{0.1111}}
  >{\raggedleft\arraybackslash}p{(\columnwidth - 12\tabcolsep) * \real{0.1481}}
  >{\raggedleft\arraybackslash}p{(\columnwidth - 12\tabcolsep) * \real{0.1481}}
  >{\raggedleft\arraybackslash}p{(\columnwidth - 12\tabcolsep) * \real{0.1481}}
  >{\raggedleft\arraybackslash}p{(\columnwidth - 12\tabcolsep) * \real{0.1481}}
  >{\raggedleft\arraybackslash}p{(\columnwidth - 12\tabcolsep) * \real{0.1481}}
  >{\raggedleft\arraybackslash}p{(\columnwidth - 12\tabcolsep) * \real{0.1481}}@{}}
\toprule\noalign{}
\begin{minipage}[b]{\linewidth}\raggedright
Current rerun backend
\end{minipage} & \begin{minipage}[b]{\linewidth}\raggedleft
Trials
\end{minipage} & \begin{minipage}[b]{\linewidth}\raggedleft
Keygen ms
\end{minipage} & \begin{minipage}[b]{\linewidth}\raggedleft
Commit ms
\end{minipage} & \begin{minipage}[b]{\linewidth}\raggedleft
Verify ms
\end{minipage} & \begin{minipage}[b]{\linewidth}\raggedleft
Open ms
\end{minipage} & \begin{minipage}[b]{\linewidth}\raggedleft
Size B
\end{minipage} \\
\midrule\noalign{}
\endhead
\bottomrule\noalign{}
\endlastfoot
Legacy ML-KEM alpha route & 20 & 0.407 & 0.627 & 0.018 & 0.658 & 1153 \\
\end{longtable}

The corresponding smoke and proof-bundle results were all successful in
the current rerun: verify \texttt{true}, opened bit \texttt{1}, round
trip \texttt{20/20}, wrong opening rejected \texttt{20/20}, tamper
detected \texttt{20/20}, and Gödel-number uniqueness
\texttt{20/20\ unique}.

\subsection{6.3 Archived keeper-backed
line}\label{archived-keeper-backed-line}

The repository also preserves an earlier keeper-backed benchmark
artifact from a previous environment. That archived line reported
30-trial timings of 0.302 ms key generation, 0.452 ms commit, 0.011 ms
verify, and 0.489 ms open at the same 1153-byte wire format. We keep
that line because it supports the paper's implementation narrative, but
we do \emph{not} relabel it as a new rerun from the present workspace.

\section{7. Wrapper portability audit}\label{wrapper-portability-audit}

A major change in this revision is that the wrapper story is written
down honestly. Instead of simply saying that keeper-ML-KEM and
libfips203 skeleton wrappers exist, the repository now includes
dedicated wrapper skeleton files and a probe script that records what
happened when they were exercised in the current environment.

\begin{longtable}[]{@{}
  >{\raggedright\arraybackslash}p{(\columnwidth - 4\tabcolsep) * \real{0.3333}}
  >{\raggedright\arraybackslash}p{(\columnwidth - 4\tabcolsep) * \real{0.3333}}
  >{\raggedright\arraybackslash}p{(\columnwidth - 4\tabcolsep) * \real{0.3333}}@{}}
\toprule\noalign{}
\begin{minipage}[b]{\linewidth}\raggedright
Wrapper probe
\end{minipage} & \begin{minipage}[b]{\linewidth}\raggedright
Current status
\end{minipage} & \begin{minipage}[b]{\linewidth}\raggedright
Observed result
\end{minipage} \\
\midrule\noalign{}
\endhead
\bottomrule\noalign{}
\endlastfoot
\texttt{keeper\_mlkem\_skeleton} & unavailable & the installed
\texttt{mlkem} package in this workspace does not export
\texttt{ML\_KEM} and \texttt{MLKEM\_768\_PARAMETERS} at the expected top
level \\
\texttt{libfips203\_skeleton} & unavailable & the Python package
\texttt{fips203} is not installed in the current workspace \\
\end{longtable}

This is a better result than vague optimism. It gives a future
maintainer two exact portability tasks: either install the Keeper-style
package that matches the documented API, or adjust the wrapper to the
actually installed module layout; and separately install the native
\texttt{libfips203} shared library in a location the Python binding can
load.

The wrapper results are captured in JSON files in
\texttt{results/wrapper\_probe\_keeper\_mlkem\_skeleton.json},
\texttt{results/wrapper\_probe\_libfips203\_skeleton.json}, and
\texttt{results/wrapper\_probe\_bundle.json}.

\section{8. Reproducibility
safeguards}\label{reproducibility-safeguards}

The repository is intended to be uploadable as a public GitHub project
rather than as an internal scratch directory. To support that use case,
the paper and repository now include the following reproducibility
safeguards.

First, the benchmark and probe results are stored as JSON rather than
only as prose. Second, the claim-guard script checks for both banned
phrases and missing disclaimer language. The script fails the build if
unsupported deployment language or off-scope positioning language is
reintroduced. Third, a minimal GitHub Actions workflow is included so
that basic checks run automatically on push. Fourth, the manuscript and
repository both keep the same sharp distinction between archived
continuity lines and current reruns.

These safeguards matter because the main failure mode of this project is
not lack of ideas; it is drift between the paper, the code, and the
measured artifacts. The present revision is explicitly designed to
reduce that drift.

\section{9. Relation to the broader research
corpus}\label{relation-to-the-broader-research-corpus}

The broader uploaded corpus includes substantial material on TICE
qubits, Floquet codes, dynamical decoupling, κ-meter metrology, and
information-curvature frameworks. Some of those materials contain
code-heavy implementations and patent-style descriptions of
quantum-control architectures, stabilizer logic, and metrology devices.
That larger corpus is valuable project context, but it is intentionally
kept out of the theorem statements and benchmark claims of this
cryptography paper.

The reason is methodological. A cryptography manuscript becomes weaker,
not stronger, if it mixes concrete implementation claims with
speculative adjacent claims about cosmology, symbolic gravity, or
generalized information-curvature physics. The current paper therefore
uses the broader corpus only in the narrow sense that it reinforced the
importance of code packaging, explicit scripts, and reproducibility
artifacts. The security claims in this paper do not depend on the TICE
or κ-meter frameworks.

\section{10. Limitations and future
work}\label{limitations-and-future-work}

This project remains a reference prototype.

We do not claim production readiness. We do not claim constant-time
behavior for the Python orchestration layer.

\begin{enumerate}
\def\labelenumi{\arabic{enumi}.}
\tightlist
\item
  \textbf{No production-readiness claim.} The code is packaged for
  public review and reproducibility, not deployment.
\item
  \textbf{No constant-time claim for Python glue.} Secret-dependent work
  should remain delegated to native or hybrid cryptographic backends;
  the Python layer should be treated as orchestration only.
\item
  \textbf{No complete computational proof yet.} The manuscript still
  does not provide a full IND-CPA, IND-CCA, or UC proof for the complete
  construction.
\item
  \textbf{Canonicalization remains part of the security surface.} If two
  semantically identical opening relations can be encoded in materially
  different ways, the neat transcript-hash argument weakens.
\item
  \textbf{Wrapper maturity is incomplete.} The keeper-style and
  libfips203 wrapper skeletons are useful, but the present environment
  probe shows that neither is turnkey in this workspace.
\item
  \textbf{Archived lines are not present reruns.} The continuity lines
  are preserved because they were explicitly requested, but they must
  not be confused with measurements produced by the current rerun
  environment.
\end{enumerate}

The next concrete steps are straightforward: complete a full
computational proof in a standard framework; stabilize one native ML-KEM
backend that works reproducibly in CI; add deterministic test vectors;
and, only after that, evaluate constant-time behavior at the compiled
binary layer rather than at the Python layer.

\section{11. Conclusion}\label{conclusion}

The cleaned-up picture is better than the earlier one. We now have a
narrower and more honest manuscript, a public repository that reflects
the paper, archived continuity lines preserved without being relabeled,
a current workspace rerun line, and explicit wrapper-portability probe
results. The cryptographic body remains ordinary and post-quantum. The
incompleteness component remains a meta-security layer for
proof-carrying adversaries only. That separation is the main
intellectual contribution of the paper, and the reproducibility tooling
is the main practical contribution.

This version should therefore be understood as a ten-page research
prototype paper with artifacts, not as a final cryptographic proof
package. That framing is more modest, but it is also much stronger.

\section{12. Artifact inventory and
workflow}\label{artifact-inventory-and-workflow}

For a public repository to be useful, a reader needs to know exactly
which script produced which artifact. This revision therefore keeps a
simple one-script-per-purpose layout.

\begin{longtable}[]{@{}
  >{\raggedright\arraybackslash}p{(\columnwidth - 4\tabcolsep) * \real{0.3333}}
  >{\raggedright\arraybackslash}p{(\columnwidth - 4\tabcolsep) * \real{0.3333}}
  >{\raggedright\arraybackslash}p{(\columnwidth - 4\tabcolsep) * \real{0.3333}}@{}}
\toprule\noalign{}
\begin{minipage}[b]{\linewidth}\raggedright
Script
\end{minipage} & \begin{minipage}[b]{\linewidth}\raggedright
Purpose
\end{minipage} & \begin{minipage}[b]{\linewidth}\raggedright
Primary output
\end{minipage} \\
\midrule\noalign{}
\endhead
\bottomrule\noalign{}
\endlastfoot
\texttt{00\_backend\_probe.py} & inspect backend metadata and
wire-format sizing &
\texttt{backend\_probe\_\textless{}backend\textgreater{}.json} \\
\texttt{01\_smoke\_test.py} & minimal end-to-end commit/verify/open
check & \texttt{smoke\_\textless{}backend\textgreater{}.json} \\
\texttt{02\_proof\_bundle.py} & round-trip and tamper tests &
\texttt{proof\_\textless{}backend\textgreater{}.json} \\
\texttt{03\_benchmark.py} & timing measurements &
\texttt{benchmark\_\textless{}backend\textgreater{}.json} \\
\texttt{04\_cross\_backend\_validate.py} & multi-backend matrix run &
\texttt{cross\_backend\_validation.json} \\
\texttt{05\_probe\_wrapper\_status.py} & keeper/libfips203
skeleton-wrapper audit & \texttt{wrapper\_probe\_bundle.json} \\
\texttt{check\_claims.py} & manuscript language guard & process exit
code and console report \\
\end{longtable}

A reviewer can therefore reproduce the repository evidence in a
predictable order: probe the selected backend, run the smoke test, run
the proof bundle, run the benchmark, then run the wrapper-portability
audit. The claim guard belongs in the same workflow because the paper is
part of the artifact surface. If the claim language drifts while the
code stays the same, the artifact package has still become less
trustworthy.

The result files are intentionally plain JSON. They are meant to be
diffable in version control, easy to upload as GitHub Actions artifacts,
and easy to quote selectively in the manuscript. This is also why the
paper includes both archived continuity lines and current rerun lines.
The goal is not to present one perfect number. The goal is to preserve
provenance.

\section{13. Threats to validity}\label{threats-to-validity}

Three threats to validity deserve explicit attention.

First, backend availability is environment-sensitive. The current
workspace can execute the legacy alpha route but not the keeper-style or
libfips203 skeleton wrappers. That does not falsify the architecture,
but it does mean that portability claims must be kept narrow. A future
reader should treat the wrapper section as a precise environment report,
not as a universal statement about those packages.

Second, microbenchmarks are fragile. Numbers in the sub-millisecond
range depend on interpreter version, CPU state, package build options,
and whether a fast path or pure-Python path was selected. For that
reason, the paper does not claim that 0.401 ms or 0.566 ms is the
uniquely correct key-generation number. It claims only that those lines
were observed in the stated artifact trail and that the 1153-byte wire
format remained stable across them.

Third, proof-carrying language can easily be misunderstood. Readers may
overread the term ``Gödel-hardness'' and assume it means a new
unconditional cryptographic hardness class. That is not what is being
claimed. The paper uses the term only for the certified-adversary game
in which the attacker is required to provide a formal proof of the bad
event. Outside that game, the scheme lives or dies by ordinary
post-quantum cryptographic reasoning.

These threats to validity are not embarrassing edge cases. They are part
of the honest shape of the current work. Recording them improves the
paper.

\section{References}\label{references}

\begin{enumerate}
\def\labelenumi{\arabic{enumi}.}
\tightlist
\item
  Kurt Gödel, \emph{Über formal unentscheidbare Sätze der Principia
  Mathematica und verwandter Systeme I}, 1931.
\item
  A. M. Turing, \emph{On Computable Numbers, with an Application to the
  Entscheidungsproblem}, 1936.
\item
  National Institute of Standards and Technology, \emph{FIPS 203:
  Module-Lattice-Based Key-Encapsulation Mechanism Standard}, 2024.
\item
  Joppe Bos et al., \emph{CRYSTALS-Kyber Algorithm Specifications and
  Supporting Documentation}, 2021.
\item
  Public repository artifacts bundled with this manuscript: benchmark
  JSONs, wrapper-probe JSONs, and claim-guard script.
\end{enumerate}

\end{document}
